\chapter{Future Work}

\section{Purpose of this Chapter}

\begin{remark}
The AIM framework presented in this manuscript establishes a
mathematical and conceptual foundation for alignment-centred modelling.
\end{remark}

\begin{remark}
However, many aspects remain open and provide opportunities for further
research, development, and practical implementation.
\end{remark}

\begin{remark}
This chapter outlines potential directions for future work, ranging from
mathematical extensions to software architecture and engineering
applications.
\end{remark}

---

\section{Extension of the Dynamic Model}

\subsection{From simplified dynamics to vehicle models}

\begin{remark}
The current formulation relies on simplified dynamic quantities such as
effective lateral acceleration and jerk.
\end{remark}

\begin{remark}
Future work should extend this toward more realistic vehicle models,
including:
\begin{itemize}
\item multi-body vehicle dynamics,
\item wheel--rail interaction,
\item suspension behaviour,
\item track elasticity.
\end{itemize}
\end{remark}

\subsection{Speed-dependent modelling}

\begin{remark}
The integration of variable speed profiles
\[
v(s)
\]
opens the possibility for more realistic evaluation and optimization of
alignments under operational conditions.
\end{remark}

---

\section{Schwerpunkt-Trassierung}

\subsection{Conceptual development}

\begin{remark}
The idea of designing the trajectory of a dynamically relevant point,
such as the center of mass, represents a fundamental extension of
classical alignment design.
\end{remark}

\begin{remark}
Further work is required to:
\begin{itemize}
\item formalize the mapping from track geometry to vehicle trajectory,
\item incorporate full dynamic models,
\item define appropriate objective functions.
\end{itemize}
\end{remark}

\subsection{Integration into optimization}

\begin{remark}
A key research direction is the coupling of Schwerpunkt-Trassierung with
optimization methods, enabling:
\begin{itemize}
\item comfort-driven design,
\item dynamic optimization,
\item direct evaluation of vehicle response.
\end{itemize}
\end{remark}

---

\section{Advanced Optimization Methods}

\subsection{SQP extensions}

\begin{remark}
While Sequential Quadratic Programming provides a powerful baseline,
future work may include:
\begin{itemize}
\item tailored SQP formulations for sparse alignment structures,
\item improved constraint handling,
\item adaptive weighting strategies.
\end{itemize}
\end{remark}

\subsection{Global optimization}

\begin{remark}
Alignment optimization problems are often non-convex.
\end{remark}

\begin{remark}
Future research may explore:
\begin{itemize}
\item global optimization techniques,
\item multi-start strategies,
\item hybrid deterministic--stochastic methods.
\end{itemize}
\end{remark}

---

\section{Network-Level Modelling}

\subsection{From single alignment to network}

\begin{remark}
Real railway systems consist of interconnected alignments forming
complex networks.
\end{remark}

\begin{remark}
Future work should extend the AIM framework to:
\begin{itemize}
\item branching structures,
\item switches and crossings,
\item corridor-wide optimization.
\end{itemize}
\end{remark}

\subsection{Topological constraints}

\begin{remark}
A systematic treatment of topological constraints remains a major
challenge.
\end{remark}

\begin{remark}
This includes:
\begin{itemize}
\item node consistency,
\item multi-track coupling,
\item network-wide feasibility conditions.
\end{itemize}
\end{remark}

---

\section{Constraint and Objective Libraries}

\subsection{Constraint formalization}

\begin{remark}
Although many railway rules can be expressed as constraints, a complete
formalization of regulatory frameworks remains an open task.
\end{remark}

\begin{remark}
Future work should aim at:
\begin{itemize}
\item systematic encoding of standards,
\item validation against real projects,
\item parameterization for different railway systems.
\end{itemize}
\end{remark}

\subsection{Objective design}

\begin{remark}
The design of objective functions is crucial for optimization-based
alignment design.
\end{remark}

\begin{remark}
Potential extensions include:
\begin{itemize}
\item cost models,
\item maintenance optimization,
\item energy efficiency,
\item lifecycle considerations.
\end{itemize}
\end{remark}

---

\section{Integration with BIM and Data Standards}

\subsection{Deeper BIM integration}

\begin{remark}
The integration of AIM with BIM workflows can be further developed by:
\begin{itemize}
\item tighter coupling with IFC alignment models,
\item bidirectional synchronization,
\item embedding AIM states within BIM objects.
\end{itemize}
\end{remark}

\subsection{Standard extensions}

\begin{remark}
The AIM framework may also inspire extensions to existing standards,
for example:
\begin{itemize}
\item inclusion of dynamic properties,
\item support for optimization metadata,
\item representation of control-based alignment models.
\end{itemize}
\end{remark}

---

\section{Interactive and Real-Time Systems}

\subsection{Live optimization}

\begin{remark}
One of the most promising directions is the development of systems that
allow real-time feedback during alignment design.
\end{remark}

\begin{remark}
This includes:
\begin{itemize}
\item live preview during optimization,
\item interactive constraint adjustment,
\item immediate visualization of dynamic effects.
\end{itemize}
\end{remark}

\subsection{User interaction}

\begin{remark}
The grabbeltisch concept may be extended to:
\begin{itemize}
\item interactive constraint definition,
\item visual debugging of alignment states,
\item intuitive editing of transition parameters.
\end{itemize}
\end{remark}

---

\section{AI and Assisted Design}

\subsection{Data-driven approaches}

\begin{remark}
Machine learning techniques may complement the AIM framework by:
\begin{itemize}
\item learning from existing alignments,
\item suggesting initial solutions,
\item identifying patterns in design choices.
\end{itemize}
\end{remark}

\subsection{Hybrid systems}

\begin{remark}
A promising direction is the combination of:
\begin{itemize}
\item physics-based modelling (AIM),
\item data-driven methods (AI),
\item interactive user control.
\end{itemize}
\end{remark}

---

\section{AXTRAN Revival and Beyond}

\begin{remark}
Classical systems such as AXTRAN demonstrated the feasibility of
optimization-based alignment design.
\end{remark}

\begin{remark}
Future work may aim to:
\begin{itemize}
\item reconstruct and generalize such systems,
\item extend them to network-level problems,
\item integrate modern numerical and computational techniques.
\end{itemize}
\end{remark}

---

\section{Software Architecture and Deployment}

\subsection{Scalability}

\begin{remark}
Practical applications require scalable implementations capable of
handling large datasets and complex networks.
\end{remark}

\subsection{Multi-view systems}

\begin{remark}
Future systems may include:
\begin{itemize}
\item synchronized multi-window views,
\item combined 2D/3D representations,
\item integration of GIS and CAD functionality.
\end{itemize}
\end{remark}

---

\section{Long-Term Vision}

\begin{remark}
In the long term, the AIM framework may contribute to a new paradigm of
infrastructure modelling in which:
\begin{itemize}
\item alignment becomes the central organizing structure,
\item geometry, dynamics, and optimization are unified,
\item engineering design becomes an interactive computational process.
\end{itemize}
\end{remark}

---

\section{Summary}

\begin{summary}
The AIM framework opens a wide range of research and development
directions.

These include extensions of the dynamic model, integration with BIM,
network-level optimization, real-time systems, and AI-assisted design.

Together, these directions point toward a future in which railway
alignment is no longer treated as static geometry, but as a dynamic,
computational, and optimizable system at the core of infrastructure
engineering.
\end{summary}
